% Source-aligned English for printed p.188.
\clearpage
\SourcePage{188}

\LayerTarget{tohoku:remark:ext-bifunctoriality:p1}
\LayerTarget{tohoku:remark:ext-center-modules}
Finally, to complete this chain of observations, note that
$\mathrm{Ext}_{\mathbf O}^p(X;A,B)$ is a module over $\Gamma({\mathbf O}^\natural)$, while
$\mathbf{Ext}_{\mathbf O}^p(A,B)$ is a module over the sheaf of rings ${\mathbf O}^\natural$, the
center of ${\mathbf O}$.

\LayerTarget{tohoku:construction:hom-composite-functors}
Formula (4.1.2) can also be written as follows. Let $h_A$ and, respectively,
$\mathbf{h}_A$ denote the covariant functors from $\mathbf C^{\mathbf O}$ to the category $\mathbf C$ of
abelian groups and, respectively, to $\mathbf C^X$, defined by
\LayerTarget{tohoku:eq:4.2.4}
\[
\tag{4.2.4}
\begin{gathered}
h_A(B)=\mathrm{Hom}_{\mathbf O}(A,B),\qquad
\mathbf{h}_A(B)=\mathbf{Hom}_{\mathbf O}(A,B):\\
h_A=\Gamma\mathbf{h}_A.
\end{gathered}
\]
Thus $h_A$ appears as a composite functor, to which Theorem 2.4.1 applies by
Proposition 4.1.3 and the corollary to Proposition 3.3.1. We therefore obtain:

\LayerTarget{tohoku:theorem:4.2.1}
\textsc{Theorem 4.2.1.} \textit{Let $X$ be a topological space
equipped with a sheaf of rings with identity ${\mathbf O}$, and let $\mathbf C^{\mathbf O}$ be the category
of ${\mathbf O}$-modules on $X$. Fix an ${\mathbf O}$-module $A\in \mathbf C^{\mathbf O}$. Then there exists on
$\mathbf C^{\mathbf O}$ a cohomological spectral functor abutting to the graded functor
$(\mathrm{Ext}_{\mathbf O}^p(X;A,B))$, whose initial term is}
\LayerTarget{tohoku:eq:4.2.5}
\[
\tag{4.2.5}
I_2^{p,q}(A,B)=H^p(X,\mathbf{Ext}_{\mathbf O}^q(A,B)).
\]

\LayerTarget{tohoku:eq:4.2.6}
In particular, one obtains the ``\textit{edge homomorphisms}''
\[
\tag{4.2.6}
\begin{aligned}
H^p(X,\mathbf{Hom}_{\mathbf O}(A,B))&\to
  \mathrm{Ext}_{\mathbf O}^p(X;A,B),\\
\mathrm{Ext}_{\mathbf O}^p(X;A,B)&\to
  H^0(X,\mathbf{Ext}_{\mathbf O}^p(A,B))
\end{aligned}
\]
(the latter is also the map arising from the interpretation of
$\mathbf{Ext}_{\mathbf O}^p(A,B)$ as the sheaf associated with the presheaf
$\mathrm{Ext}_{\mathbf O}^p(-;A,B)$), together with the five-term exact sequence
\LayerTarget{tohoku:eq:4.2.7}
\[
\tag{4.2.7}
\begin{aligned}
0\to H^1(X,\mathbf{Hom}_{\mathbf O}(A,B))
&\to \mathrm{Ext}_{\mathbf O}^1(X;A,B)
\to H^0(X,\mathbf{Ext}_{\mathbf O}^1(A,B))\\
&\to H^2(X,\mathbf{Hom}_{\mathbf O}(A,B))
\to \mathrm{Ext}_{\mathbf O}^2(X;A,B).
\end{aligned}
\]
\LayerTarget{tohoku:remark:ext1-extension-classes}
This exact sequence \textit{makes precise, in particular, the structure of
the group $\mathrm{Ext}_{\mathbf O}^1(X;A,B)$ of classes of ${\mathbf O}$-module extensions of
$A$ \textup{(the quotient module)} by $B$ \textup{(the submodule).}}

\LayerTarget{tohoku:construction:stalk-ext-comparison}
To use the spectral sequence just obtained, it remains to specify the
calculation of $\mathbf{Ext}_{\mathbf O}^p(A,B)$. Since the functor
$F\to F(x)$, which assigns to an abelian sheaf $F$ its stalk $F(x)$,
is \textit{exact}, the groups $\mathbf{Ext}_{\mathbf O}^p(A,B)(x)$ may, with $A$ fixed and $p$
and $B$ variable, be regarded as forming a universal covariant cohomological
functor in $B$, reducing in dimension 0 to $\mathbf{Hom}_{\mathbf O}(A,B)(x)$. Likewise,
the groups $\mathrm{Ext}_{{\mathbf O}(x)}^p(A(x),B(x))$ form a cohomological functor in
$B$, reducing in dimension 0 to
$\mathrm{Hom}_{{\mathbf O}(x)}(A(x),B(x))$. The functorial homomorphism (4.1.3) and the
definition of universal cohomological functors therefore give homomorphisms
\LayerTarget{tohoku:eq:4.2.8}
\[
\tag{4.2.8}
\mathbf{Ext}_{\mathbf O}^p(A,B)(x)\to
\mathrm{Ext}_{{\mathbf O}(x)}^p(A(x),B(x))
\]
(characterized by defining a homomorphism of cohomological functors which in
dimension 0 reduces to homomorphism (4.1.3)).

\TohokuContinuation{189}
